% main.tex -- arXiv v1 skeleton for research/paper/ (gaps C1/C2 in
% research/05-arxiv-readiness.md).
%
% Paper:
%   "How Fragile Are Deployed Text Watermarks? An Empirical Study of
%    Layered Watermark Removal under Realistic User-Side Editing"
%
% Sources of truth: research/02-paper-outline.md (title 01/02 §1,
% abstract §2, section skeleton §5, tables/figures spec §6, citation
% map §7), research/03-related-work.md (verified bibliography).
%
% Status: SKELETON ONLY. Placeholder prose + TODO comments; NO real
% experiment numbers yet (real values land after the W1 core run; see
% research/01-experiment-protocol.md and research/05-arxiv-readiness.md
% gaps A/B/C).
%
% Compile (no ACL style package required for the skeleton):
%   pdflatex main
%   bibtex  main
%   pdflatex main
%   pdflatex main
%
% At submission: drop in acl2024.sty and switch the preamble to the ACL
% template; the \input{} files and the body carry over. Also swap
% \bibliographystyle{plainnat} for the ACL bib style if required.
%
% NOTE: tables/ and figures/ do not exist yet -- they are generated by
% research/scripts/make_tables.py and research/scripts/make_figures.py
% (gap C3). main.tex will not compile until those files are emitted.

\documentclass[10pt,letterpaper]{article}

% --- packages ------------------------------------------------------------
% Standard packages only; the skeleton deliberately avoids venue-specific
% styles (ACL style is dropped in at submission, see the header comment).
\usepackage[T1]{fontenc}   % UTF-8 is the default for modern engines
\usepackage{graphicx}      % \includegraphics{figures/fN} (make_figures.py)
\usepackage{booktabs}      % tables generated by make_tables.py
\usepackage{natbib}        % \citet / \citep (author-year by default)
\usepackage[hidelinks]{hyperref}

% --- metadata ------------------------------------------------------------
% Title locked in research/02-paper-outline.md §1.
\title{How Fragile Are Deployed Text Watermarks? An Empirical Study of
Layered Watermark Removal under Realistic User-Side Editing}

% Solo author, no affiliation (allowed on arXiv; see
% research/05-arxiv-readiness.md E1). Affiliation is optional at
% submission -- uncomment and fill in if desired, e.g.:
%   \author{Guillaume Meyer\\ \small{Affiliation}}
\author{Guillaume Meyer}
\date{}   % arXiv stamps its own submission date.

\begin{document}

\maketitle

% --- abstract ------------------------------------------------------------
% Draft from research/02-paper-outline.md §2 lives in abstract.tex.
% TODO: recheck all numbers against the final run (writing phase W2).
\begin{abstract}
% abstract.tex -- abstract draft for the arXiv v1 paper.
% Source: research/02-paper-outline.md §2 (v0.2, 2026-08-18), copied
% verbatim (LaTeX-escaped). ~150 words.
%
% TODO (W2): recheck every number against the final run before
% submission -- the 3,500/3,500 cell counts and the metric list must
% match the released results (research/01-experiment-protocol.md §2;
% 05-arxiv-readiness.md checklist F).

Text watermarking is the primary mechanism proposed for EU AI Act
Art.~50 transparency obligations on machine-generated content. We
measure its robustness against realistic user-side editing. Using
same-config detection over 3,500 watermarked and 3,500 unwatermarked
generations (KGW, SynthID-Text, EXP, Unigram, and SIR schemes, in
English and German/French/Spanish), we evaluate a layered removal
pipeline that combines formatting-layer cleanup (invisible Unicode,
bidi) with statistical rewriting driven by detection feedback, and we
report ROC-based metrics (AUROC, TPR@1\%FPR) and quality metrics
(perplexity, BERTScore, ROUGE-L) at the point of detection collapse. We
find that quality-preserving paraphrase and translation round-trip
collapse KGW-class detection to near chance, that SynthID-Text resists
token substitution but not paraphrase, that layering dominates single
attacks at equal quality cost, and that multilingual texts are
systematically more fragile. We discuss implications for Art.~50
compliance and release our corpus, configs, and harness.

\end{abstract}

% =====================================================================
% 1. Introduction
% =====================================================================
\section{Introduction}
% Hook (02 §5.1): Art. 50 in force 2026-08-02; vendor rollouts (Claude,
% Gemini); Google retired the SynthID-Text API in Aug 2026 -- the market
% is consolidating on methods that do not survive editing. TODO: one
% sentence of this market context here, no metrics needed.

The EU AI Act's transparency obligations for machine-generated content
(Regulation (EU) 2024/1689, Art.~50) \citep{euaiact2024} entered force
on 2~August 2026, and text watermarking is the primary mechanism
proposed to meet them. % TODO: add the vendor-rollout / SynthID-Text
% retirement context (one sentence).

% Motivation (02 §5.1): the public deployment of our own removal tool --
% users edit their own output; do watermarks survive? (1 short
% paragraph, no metrics needed.)
The public deployment of a watermark-removal tool in this project showed
that ordinary users routinely edit text they generated themselves
before publishing it. That raises the basic question we study: do
deployed text watermarks survive realistic user-side editing?

% Contributions (02 §4) -- TODO: verify the four bullets at submission.
Our contributions are as follows:
\begin{itemize}
  \item \textbf{Measurement.} The first ROC-based robustness study of
  deployed-class text watermarking (KGW, SynthID-Text, EXP, Unigram,
  SIR) under realistic, layered user-side editing, with empirical-null
  FPR calibration.
  \item \textbf{Method.} A layered removal pipeline (formatting +
  statistical) with detection-feedback rewriting; we show it dominates
  single-layer baselines at equal quality cost (Pareto).
  \item \textbf{Resource.} An open corpus (25 EN + 75 multilingual
  prompts), a config-pinned harness (MarkLLM-based), and results
  (JSONL) for reproducible attack/defense benchmarking.
  \item \textbf{Policy measurement.} Evidence on whether Art.~50
  transparency obligations can rely on watermarking, with concrete
  recommendations (platform-level provenance, metadata,
  robust-but-invisible schemes, honest failure modes).
\end{itemize}

Figure~\ref{fig:f1} illustrates the layered pipeline we evaluate;
Section~\ref{sec:background} situates the work relative to the
watermarking literature.

% =====================================================================
% 2. Background and Related Work
% =====================================================================
\section{Background and Related Work}
\label{sec:background}
% Watermarking families (02 §5.2; citation map 02 §7; full list 03 A).
Sampling-based watermarking biases the token distribution of the
generator under a secret key, starting with KGW
\citep{kirchenbauer2023kwg} and a large family of successors
\citep{christ2023undetectable,kuditipudi2023robust,hu2023unbiased,
wu2023dipmark,zhao2024permute,huo2024tswatermark,liu2024adaptive}.
Tournament-based methods such as SynthID-Text
\citep{deepmind2024synthidtext,omidi2026synthid} pair tokens into
tournaments to avoid distribution distortion. Semantic schemes
watermark in embedding space \citep{liu2023sir,hou2023semstamp,
hou2024ksemstamp}, provable families aim for statistical guarantees
\citep{liu2023upv}, and recent work adds entropy-aware and adaptive
designs \citep{gu2025invisible,wang2025morphmark} as well as
black-box watermarking \citep{yang2023blackbox}. We follow the
taxonomy of \citet{liu2023survey} and build on the MarkLLM toolkit
\citep{pan2024markllm}.

% Attack literature (02 §5.2; citation map 02 §7; full list 03 B/C).
Robustness studies established that watermarks are fragile under
paraphrasing and other edits \citep{kirchenbauer2023reliability}, that
strong watermarking faces impossibility limits under random-walk
attacks \citep{zhang2023sand,harelcanada2025sandcastles}, and that
translation round-trips defeat several schemes \citep{he2024xsir};
watermark keys can also be stolen from a black-box detector
\citep{jovanovic2024stealing}. Closest to us, recent work assesses
SynthID-Text robustness specifically \citep{han2025synthid,
omidi2026synthid} and the forensic readiness of watermark evidence
\citep{tamim2026forensic}. Detector-side methods improve segment
detection \citep{pan2024waterseeker,lu2024ewd}, users can sometimes
identify watermarked outputs through crafted prompts
\citep{liu2024crafted}, and watermarks may not robustly prevent
knowledge distillation \citep{pan2025distillation}. Metadata
provenance standards such as C2PA \citep{c2pa2024spec} are
complementary; a metadata mini-study is deferred to v2.

% Positioning (02 §5.2; 03 "Positioning summary").
Relative to this literature, we add (i) a \emph{layered} attack that
combines a formatting layer with statistical rewriting, (ii) a
systematic ROC measurement of deployed-class schemes under a single
protocol, (iii) a quality--detectability Pareto analysis, and (iv) an
Art.~50 policy measurement; none of the closest neighbors
\citep{he2024xsir,han2025synthid,omidi2026synthid,tamim2026forensic,
kirchenbauer2023reliability} combines all four.

% =====================================================================
% 3. Threat Model and System
% =====================================================================
\section{Threat Model and System}
% Who (02 §5.3): end users editing text they generated with their own
% account -- explicitly NOT third-party content (see ethics.tex).
We consider users who edit text they generated themselves with their
own account before publishing it, for reasons of privacy or style. We
do not attack third-party content, and we do not frame the study as an
attacker--victim setting.

% What (02 §5.3): watermark-as-label (not access control); detector =
% same-config MarkLLM detector (standard in the literature; vendor
% detectors are black-box and out of scope -- Section 8).
Watermarking is a label, not a lock: no access control is bypassed. The
detector is the same-config open-source MarkLLM implementation
\citep{pan2024markllm} used for generation, following standard
practice in the robustness literature; vendor detectors are not probed
(Section~\ref{sec:limitations}).

% System (02 §5.3): two layers for v1 -- A formatting / B statistical --
% with a detection-feedback rewrite loop (Fig. 1).
Our v1 pipeline has two layers: Layer~A strips formatting-layer marks
(invisible Unicode, bidi, tag cleanup) deterministically, and Layer~B
applies statistical rewriting (paraphrase, back-translation, structural
rewrite, humanization) driven by a detection-feedback loop that stops
once detection collapses. Figure~\ref{fig:f1} shows the system.

\begin{figure}[t]
  \centering
  \includegraphics[width=\linewidth]{figures/f1}
  \caption{The layered removal pipeline: Layer~A (formatting cleanup)
  followed by Layer~B (statistical rewriting with a detection-feedback
  loop).}
  \label{fig:f1}
\end{figure}
% TODO: figures/f1.pdf is generated by research/scripts/make_figures.py
% (pipeline diagram; gap C3) -- until it exists the skeleton will not
% compile.

% Definitions box (02 §5.3).
\begin{quote}
\textbf{Definitions.} \emph{TPR@FPR} -- true-positive rate at a fixed
false-positive rate (we report TPR@1\%FPR); \emph{AUROC} -- area
under the ROC curve; \emph{detection collapse} -- detection at
near-chance levels (AUROC $\approx 0.5$); \emph{quality cost} --
drift in perplexity, BERTScore, and ROUGE-L relative to the unedited
output. % TODO: freeze exact metric definitions against 01 §5 once B2
% lands.
\end{quote}

% =====================================================================
% 4. Experimental Setup
% =====================================================================
\section{Experimental Setup}
% Design summary (02 §5.4 -> 01 §2): locked v1 factorial.
We follow the locked v1 factorial design of
\emph{research/01-experiment-protocol.md}: seven scheme configurations
(KGW $\gamma{=}.25,\delta{=}1$; KGW $\gamma{=}.5,\delta{=}2$; KGW
$\gamma{=}.5,\delta{=}4$; SynthID-Text; EXP; Unigram; SIR) across
lengths $\{100,300,500\}$, temperatures $\{0.7,1.0\}$, and
languages en, de, fr, es (restricted subsets) -- 3,500 watermarked and
3,500 unwatermarked generations (7,000 texts).
% TODO: cross-check the cell counts against the final run (01 §2).

% Generation protocol (01 §3).
English texts are generated with \texttt{facebook/opt-1.3b} via
MarkLLM; the multilingual cells use \texttt{Qwen/Qwen2.5-1.5B-Instruct}
as an explicit model holdout. Watermark configs are pinned in
\texttt{research/configs/} and the same JSON is used for generation
and detection; seeds are fixed and recorded.

% Attack cells A0-A8 (01 §4) + T1 taxonomy table (02 §6).
Each text passes through the attack cells A0--A8 of the protocol
(Table~\ref{tab:t1}): no attack, Layer~A only, single-pass paraphrase,
adaptive paraphrase with early stopping, back-translation round-trip,
structural rewrite, humanization, cheap token-level baselines, and the
full layered pipeline (A1 then A2/A3).

% Detection + quality (01 §5): empirical null; quality metrics.
Detection uses the same-config MarkLLM detectors with an empirical null
distribution estimated from unwatermarked controls rather than normal
assumptions; quality is measured with perplexity (gpt2-large),
BERTScore, ROUGE-L, and length drift on a stratified subset.

% Reproducibility (01 §7-8; 05 A7/E3).
We release the corpus, configs, seeds, harness pins (MarkLLM commit, HF
revisions, pip freeze), and results as JSONL.
% TODO: manifest.json + release URLs before submission (gaps A7/E3).

\begin{table}[t]
  \centering
  \caption{Attack taxonomy: family $\times$ mechanism $\times$
  implementation $\times$ prior-art anchor.}
  \label{tab:t1}
  \input{tables/t1}
\end{table}
% TODO: tables/t1.tex is generated by research/scripts/make_tables.py
% (real rows from the run; gap C3).

% =====================================================================
% 5. Results
% =====================================================================
\section{Results}
% Money table (02 §6 T2).
Table~\ref{tab:t2} reports the core AUROC / TPR@1\%FPR matrix: rows
are scheme configurations (7) $\times$ attacks (8); columns are
pre-attack, A-only, B-only, and A+B.
% TODO: fill Table 2 with real numbers from the run via
% research/scripts/make_tables.py, and confirm the five key claims
% (02 §5.5):
%   (i)   paraphrase/back-translation collapse KGW-class detection;
%   (ii)  SynthID-Text resists token substitution but not paraphrase;
%   (iii) layering > single layers at equal quality cost;
%   (iv)  even the strongest config (gamma=.5, delta=4, 300 tok) drops
%         below usable TPR under adaptive rewrite;
%   (v)   multilingual (de/fr/es) is systematically more fragile.

Figure~\ref{fig:f2} shows ROC curves before and after removal per
scheme, and Figure~\ref{fig:f3} shows the quality--detectability
Pareto frontier. Table~\ref{tab:t3} reports quality metrics
(perplexity delta, BERTScore, ROUGE-L, length drift, number/URL
survival) at the point of detection collapse; Table~\ref{tab:t4} and
Figure~\ref{fig:f4} ablate watermark strength ($\gamma,\delta$) by
length and temperature; Table~\ref{tab:t5} covers the multilingual
grid; Table~\ref{tab:t6} reports attack cost (tokens in/out, wall
time, USD per 1k words); Table~\ref{tab:t7} compares with published
baselines where numbers are reported in the same metric.
Figure~\ref{fig:f5} shows a redacted before/after case study with
detector scores.
% TODO: replace the placeholder prose above with result-specific
% discussion once the run lands (writing phase W2).

\begin{table}[t]
  \centering
  \caption{AUROC / TPR@1\%FPR matrix (the money table).}
  \label{tab:t2}
  \input{tables/t2}
\end{table}
% TODO: tables/t2.tex from research/scripts/make_tables.py.

\begin{figure}[t]
  \centering
  \includegraphics[width=\linewidth]{figures/f2}
  \caption{ROC curves pre/post removal per scheme (the money figure).}
  \label{fig:f2}
\end{figure}
% TODO: figures/f2.pdf from research/scripts/make_figures.py.

\begin{figure}[t]
  \centering
  \includegraphics[width=\linewidth]{figures/f3}
  \caption{Quality--detectability Pareto frontier (PPL vs.\ AUROC
  across attacks).}
  \label{fig:f3}
\end{figure}
% TODO: figures/f3.pdf from research/scripts/make_figures.py.

\begin{table}[t]
  \centering
  \caption{Quality per attack at the collapse point: PPL delta,
  BERTScore, ROUGE-L, length drift, number/URL survival.}
  \label{tab:t3}
  \input{tables/t3}
\end{table}
% TODO: tables/t3.tex from research/scripts/make_tables.py.

\begin{table}[t]
  \centering
  \caption{Ablation: TPR@1\%FPR $\times$ ($\gamma,\delta$) $\times$
  length $\times$ temperature.}
  \label{tab:t4}
  \input{tables/t4}
\end{table}
% TODO: tables/t4.tex from research/scripts/make_tables.py.

\begin{figure}[t]
  \centering
  \includegraphics[width=\linewidth]{figures/f4}
  \caption{TPR@1\%FPR vs.\ watermark strength, length as line
  style.}
  \label{fig:f4}
\end{figure}
% TODO: figures/f4.pdf from research/scripts/make_figures.py.

\begin{table}[t]
  \centering
  \caption{Multilingual results (en/de/fr/es; v1 extension).}
  \label{tab:t5}
  \input{tables/t5}
\end{table}
% TODO: tables/t5.tex from research/scripts/make_tables.py.

\begin{table}[t]
  \centering
  \caption{Attack cost: tokens in/out, wall time, USD per 1k words.}
  \label{tab:t6}
  \input{tables/t6}
\end{table}
% TODO: tables/t6.tex from research/scripts/make_tables.py.

\begin{table}[t]
  \centering
  \caption{Comparison with published baselines (cite-and-compare where
  numbers are reported in the same metric).}
  \label{tab:t7}
  \input{tables/t7}
\end{table}
% TODO: tables/t7.tex from research/scripts/make_tables.py.

\begin{figure}[t]
  \centering
  \includegraphics[width=\linewidth]{figures/f5}
  \caption{Case study: redacted before/after text with detector
  scores.}
  \label{fig:f5}
\end{figure}
% TODO: figures/f5.pdf from research/scripts/make_figures.py.

% =====================================================================
% 6. Analysis and Case Study
% =====================================================================
\section{Analysis and Case Study}
% Cost of attack vs cost of defense (02 §5.6; Table 6).
The cost asymmetry between attack and defense (Table~\ref{tab:t6}) is
central to the policy argument: an attacker spends minutes and a few
cents of API credit per document, while the defender must raise
watermark strength at a measurable quality cost.
% TODO: quote the real cost figures and quality deltas from the run.

% Failure modes ranked (02 §5.6).
Ranking the failure modes, token-level statistical watermarks (the KGW
line) collapse under semantic rewriting, n-gram/tournament methods
(SynthID-Text) survive token substitution but not paraphrase, and
semantic-invariant properties survive best. % TODO: replace with the
% ranked, number-backed analysis once results exist; the redacted case
% study appears in Figure~\ref{fig:f5}.

% =====================================================================
% 7. Policy Discussion
% =====================================================================
\section{Policy Discussion}
% Art. 50 mechanics (04 §2; 02 §5.7; citation map 02 §7).
Article~50 of Regulation (EU) 2024/1689 \citep{euaiact2024} places
transparency obligations on providers and deployers of general-purpose
AI systems, not on end users editing their own documents: the
obligation is that machine-readable output be detectable, and our
measurements ask whether that detection survives ordinary editing.
% TODO: verify exact recital numbers before submission (04 §2).

Recent forensic-readiness evidence \citep{tamim2026forensic} is
consistent with our findings. We argue that platform-level provenance
and metadata (C2PA) \citep{c2pa2024spec} are the robust complement to
text-layer watermarking, and that regulators should not rely on
token-level schemes alone. Figure~\ref{fig:f6} places the measured
collapse against the Art.~50 timeline. % TODO: honest limits of the
% study here (no vendor black-box measurement; see Section 8).

\begin{figure}[t]
  \centering
  \includegraphics[width=\linewidth]{figures/f6}
  \caption{Art.~50 timeline vs.\ measured detection collapse (policy
  figure).}
  \label{fig:f6}
\end{figure}
% TODO: figures/f6.pdf from research/scripts/make_figures.py.

% =====================================================================
% 8. Limitations
% =====================================================================
\section{Limitations}
\label{sec:limitations}
% 02 §5.8.
Our detector is same-config and open-source; vendor detectors are
black-box and were not probed. Generators are small open-weight models
(opt-1.3b; Qwen2.5-1.5B-Instruct for de/fr/es as an explicit model
holdout) rather than frontier LLMs. API-probe limits bound the rewrite
loop, and our detection-feedback oracle is stronger than a naive user;
we therefore report both adaptive and single-pass numbers.
% TODO: add any additional limitations surfaced during the run.

% =====================================================================
% 9. Ethics Statement
% =====================================================================
\section{Ethics Statement}
% Draft from research/04-ethics-and-legal.md §3, adapted to a paragraph;
% adapt to the venue template at submission (arXiv has no ethics
% checklist -- 04 §6).
% ethics.tex -- ethics statement for the arXiv v1 paper.
% Source: research/04-ethics-and-legal.md §3 (draft, adapted to a
% paragraph) and §5 (disclosure). Keep this file in sync with 04; adapt
% to the venue template at submission (arXiv has no ethics checklist --
% 04 §6). Do not weaken the "no third-party content" and
% "watermark-as-label" framing.

All text used in this study was generated by the authors using
open-weight models (opt-1.3b) via the MarkLLM toolkit
\citep{pan2024markllm}; no third-party content, user data, or live
vendor outputs were used, and no watermarked content produced by
commercial providers was collected or altered. The removal pipeline
evaluated here operates on text generated by the same user who owns it;
the study does not enable or endorse alteration of third-party content.
Watermarking is a label, not an access-control mechanism, so no
security control is circumvented. We disclose our findings to support
(i) realistic expectations for regulators relying on watermarking for
transparency obligations (Art.~50, Regulation (EU) 2024/1689
\citep{euaiact2024}), and (ii) the design of more robust provenance
mechanisms. Detectors are same-config open-source implementations;
commercial detectors were not probed. We do not provide live removal
services or weights tuned for evasion of specific vendor detectors
beyond what is reported. The authors' tooling is public
(\url{https://github.com/guillaumemeyer/watermarks-remover}) and was
deployed publicly before this study began; this paper formalizes
measurements of mechanisms already in production use.

% Disclosure (04 §5): the tool is already public and widely reported;
% no embargo or coordinated-disclosure obligation applies to this
% measurement -- state this if a venue asks about disclosure.


% =====================================================================
% 10. Conclusion
% =====================================================================
\section{Conclusion}
We measured the robustness of deployed-class text watermarking under
realistic, layered user-side editing. Our results indicate that
quality-preserving rewriting collapses KGW-class detection, that
layering dominates single attacks at equal quality cost, and that
multilingual text is systematically more fragile. % TODO: summarize
% the final numbers and the Art. 50 recommendation here (02 §5.10).

% --- acknowledgments -----------------------------------------------------
\section*{Acknowledgments}
% Source: research/04-ethics-and-legal.md §5 + research/03-related-work.md §D.
% acknowledgments.tex -- acknowledgments for the arXiv v1 paper.
% Source: research/04-ethics-and-legal.md §5 (MarkLLM is Apache-2.0;
% attribution required) and research/03-related-work.md §D (MarkLLM:
% Pan et al., EMNLP 2024 Demo).

We thank the Tsinghua NLP group (THU-BPM) for releasing MarkLLM
\citep{pan2024markllm} -- an open-source, Apache-2.0 toolkit for LLM
watermarking -- via the Generative-Watermark-Toolkits repository, which
served as the generation and detection harness for this study. We also
thank the Hugging Face community for the open models used in our
experiments (opt-1.3b and Qwen2.5-1.5B-Instruct).
% TODO: add any further acknowledgments before submission.


% --- bibliography ---------------------------------------------------------
\bibliographystyle{plainnat}
\bibliography{refs}
% TODO: every entry in refs.bib must be cited (or pruned) before
% submission; re-verify all IDs and venue labels (03 header).

\end{document}
